\appendix

\section{The second-order independence sandwich}\label{app:second-order}

Setting of \cref{thm:janson-general}: increasing events $B_1,\dots,B_m$ (after
flipping the variables that occur positively, \cref{sec:independence}) with syntactic
supports $S_i$, a product distribution, and a fixed order. For each $i$ let
$N_i=\{j<i:j\sim i\}$, $R_i=\{j<i:j\not\sim i\}$, $\Gamma(j)$ the neighbourhood of $j$,
$b_i=D(B_i)$ with $\max_ib_i<1$, $\beta_i=\sum_{j\in N_i}b_j$ (assumed $<1$),
\[
C_i=\sum_{j\in N_i}\mathrm{Cov}(B_i,B_j)\ \ (\ge0\text{ by Harris}),
\]
\[
\Theta_i=b_i\Big[\sum_{j\in N_i}\sum_{k\in R_i\cap\Gamma(j)}D(B_jB_k)+\sum_{j<j'\in N_i}D(B_jB_{j'})\Big],
\]
\[
\Theta'_i=\sum_{j\in N_i}\sum_{k\in R_i\cap\Gamma(j)}D(B_iB_jB_k)+\sum_{j<j'\in N_i}D(B_iB_jB_{j'}).
\]
Each triple $\{i,j,k\}$ with $j,k<i$ that is connected in $G$ occurs exactly once in
these sums (in the first if $i$ is adjacent to only one of $j,k$, in the second if it
is adjacent to both), so $\Theta'_i=\sum D(B_iB_jB_k)$ and $\Theta_i=b_i\sum D(B_jB_k)$
over these triples, and $\sum_i\Theta'_i$ does not depend on the order. If $b_i=0$
then $c_i$, $C_i$, $\Theta_i$ and $\Theta'_i$ vanish and the $i$-th factors below equal $1$;
so we assume $b_i>0$ whenever we condition on $B_i$.

\begin{theorem}\label{thm:second-order}
\[
\prod_i\Big(1-\min\Big\{b_i,\;b_i-C_i+\frac{\Theta'_i}{1-\beta_i}\Big\}\Big)\;\le\;p\;\le\;
\prod_i\Big(1-\max\Big\{0,\;b_i-\frac{C_i+\Theta_i}{1-\beta_i}\Big\}\Big).
\]
Consequently, with $M_2=\prod_i(1-b_i+C_i)$,
$|\log p-\log M_2|\le E_3:=\sum_i\frac{\Theta_i+\Theta'_i+\beta_iC_i}{(1-\beta_i)(1-b_i)}$.
\end{theorem}

\begin{proof}
Fix $i$, write $N=N_i$, $R=R_i$, $F_X=\bigwedge_{j\in X}\neg B_j$ and
$\nu=D(\cdot\mid F_R)$ (by Harris $\Pr[F_R]\ge\prod_{j\in R}(1-b_j)>0$). Then $p=\prod_i(1-c_i)$ with $c_i=\Pr[B_i\mid F_N\wedge F_R]$.
Since $S_i$ is disjoint from the supports of the events in $R$, $\nu(B_i)=b_i$ and
\[
c_i=b_i-\frac{b_i\,[\nu(F_N)-\nu(F_N\mid B_i)]}{\nu(F_N)},\qquad
\nu(F_N)-\nu(F_N\mid B_i)=\nu\big(\textstyle\bigcup_NB_j\mid B_i\big)-\nu\big(\textstyle\bigcup_NB_j\big).
\]
We use six facts. Each follows from the Harris inequality (an increasing and a
decreasing event are negatively correlated under a product measure) and from
independence of events with disjoint supports. For $j\in N$ let $R'=R\cap\Gamma(j)$
and $R''=R\setminus R'$; the events of $R''$ have supports disjoint from
$S_i\cup S_j$, and $\Pr[F_R]\le\Pr[F_{R''}]$.
\begin{enumerate}[label=(H\arabic*)]
\item $\nu(B_j\mid B_i)\le D(B_iB_j)/b_i$ (Harris for $B_iB_j$ versus $F_R$).
\item $\nu(B_j)\le b_j$ (Harris).
\item $\nu(B_j)\ge b_j-\sum_{k\in R'}D(B_jB_k)$: $\Pr[B_jF_R]\ge\Pr[B_jF_{R''}]-\sum_{k\in
R'}\Pr[B_jB_kF_{R''}]\ge(b_j-\sum_kD(B_jB_k))\Pr[F_{R''}]$ by independence and Harris;
divide by $\Pr[F_R]\le\Pr[F_{R''}]$ (trivial if the bracket is negative).
\item $\nu(B_jB_{j'})\le D(B_jB_{j'})$ (Harris).
\item $\nu(B_j\mid B_i)\ge\big[D(B_iB_j)-\sum_{k\in R'}D(B_iB_jB_k)\big]/b_i$: as (H3), with
$B_iB_j$ in place of $B_j$.
\item $\nu(B_jB_{j'}\mid B_i)\le D(B_iB_jB_{j'})/b_i$ (Harris).
\end{enumerate}
\emph{Upper bound on $p$.} By Bonferroni, $\nu(\bigcup B_j\mid B_i)\le\sum_j\nu(B_j\mid
B_i)$ and $\nu(\bigcup B_j)\ge\sum_j\nu(B_j)-\sum_{j<j'}\nu(B_jB_{j'})$. With (H1), (H3),
(H4), $b_i[\nu(F_N)-\nu(F_N\mid B_i)]\le C_i+\Theta_i$; with (H2), $\nu(F_N)\ge1-\beta_i$.
So $c_i\ge b_i-(C_i+\Theta_i)/(1-\beta_i)$, and $c_i\ge0$.

\emph{Lower bound on $p$.} By Bonferroni the other way with (H5), (H6), (H2),
$b_i[\nu(F_N)-\nu(F_N\mid B_i)]\ge C_i-\Theta'_i$. If $C_i\ge\Theta'_i$ use
$\nu(F_N)\le1$; otherwise use $\nu(F_N)\ge1-\beta_i$. In both cases (using $C_i\ge0$)
$c_i\le b_i-C_i+\Theta'_i/(1-\beta_i)$; and $c_i\le b_i$ by Harris.

\emph{The $E_3$ bound.} Put $x_i=1-b_i+C_i\ge1-b_i$. On the upper side, the factor is
at most $x_i+(\Theta_i+\beta_iC_i)/(1-\beta_i)$ (or $1$, in which case
$1-x_i<(\Theta_i+\beta_iC_i)/(1-\beta_i)$), and $\log(x+y)-\log x\le y/x$. On the lower
side, the factor is at least $\max\{1-b_i,\ x_i-\Theta'_i/(1-\beta_i)\}$; if the
second term is the larger one, $\log(x-y)-\log x\ge-y/(x-y)\ge-y/(1-b_i)$, and
otherwise $C_i<\Theta'_i/(1-\beta_i)$ and $\log(1-b_i)-\log x_i\ge-C_i/(1-b_i)$, which is
no smaller. Summing over $i$ gives the claim.
\end{proof}

We checked the theorem against exact enumeration on 382 random instances (3--8
random monotone events, each a monotone DNF on 2--4 of 8--14 variables, random product
distributions and orders; \texttt{experiments/check\_second\_order.py}): no
violations, median $E_3=0.064$. The $\min$ in the lower bound matters: without it a
factor can become negative, and a product of two negative factors can exceed $p$.
